% =====================================================================
\chapter{La descomposición en valores singulares}
\label{tema:4}
% =====================================================================

El teorema espectral solo se aplica a matrices cuadradas simétricas. La descomposición en valores singulares (SVD) extiende la misma idea a cualquier matriz rectangular: toda transformación lineal es un giro, seguido de un estiramiento a lo largo de ejes perpendiculares, seguido de otro giro. Este tema sigue a Strang~\cite[§7.1]{strang-ila}, la clase~29 de 18.06~\cite{mit1806}, las clases~6 y~8 de 18.065~\cite{mit18065} y el §4.5 de~\cite{mml}.

\section{El teorema de la SVD}

\begin{teorema}[Descomposición en valores singulares]\label{teo:svd}
Sea $A\in\R^{m\times n}$ de rango $r$. Existen matrices ortogonales $U\in\R^{m\times m}$ y $V\in\R^{n\times n}$ y números $\sigma_1\ge\sigma_2\ge\dots\ge\sigma_r>0$ tales que
\[
A=U\Sigma V\T=\sum_{i=1}^{r}\sigma_i\,u_iv_i\T,
\]
donde $\Sigma\in\R^{m\times n}$ tiene $\Sigma_{ii}=\sigma_i$ para $i\le r$ y ceros en el resto, y $u_i$, $v_i$ son las columnas de $U$ y $V$. Los $\sigma_i$ son los \emph{valores singulares} de $A$, y los $u_i$ y $v_i$, sus \emph{vectores singulares} izquierdos y derechos. Además, $Av_i=\sigma_iu_i$ y $A\T u_i=\sigma_iv_i$ para $i\le r$.
\end{teorema}

\begin{proof}
Por la \cref{prop:AtA} y el \cref{teo:espectral}, $A\T A\in\R^{n\times n}$ tiene una base ortonormal de autovectores $v_1,\dots,v_n$ con autovalores $\lambda_1\ge\dots\ge\lambda_n\ge 0$. Como $\Nuc(A\T A)=\Nuc(A)$ (si $A\T Ax=0$, entonces $\norm{Ax}^{2}=x\T A\T Ax=0$), el \cref{teo:rango-nulidad} da $\rango(A\T A)=\rango(A)=r$, de modo que exactamente $\lambda_1,\dots,\lambda_r$ son positivos. Se define $\sigma_i=\sqrt{\lambda_i}$ y, para $i\le r$,
\[
u_i=\frac{Av_i}{\sigma_i}.
\]
Estos vectores son ortonormales:
\[
u_i\T u_j=\frac{v_i\T A\T Av_j}{\sigma_i\sigma_j}=\frac{\lambda_j\,v_i\T v_j}{\sigma_i\sigma_j}=\delta_{ij}.
\]
Se completan $u_1,\dots,u_r$ con $u_{r+1},\dots,u_m$ hasta una base ortonormal de $\R^{m}$ (Gram--Schmidt). Para $i>r$, $\norm{Av_i}^{2}=\lambda_i=0$, luego $Av_i=0$. Por tanto, $Av_i=\sigma_iu_i$ para $i\le r$ y $Av_i=0$ para $i>r$, que en forma matricial es $AV=U\Sigma$. Multiplicando por $V\T$ a la derecha, $A=U\Sigma V\T$. Por último, $A\T u_i=A\T Av_i/\sigma_i=\lambda_iv_i/\sigma_i=\sigma_iv_i$.
\end{proof}

\begin{figure}[h]
\centering
\begin{tikzpicture}[scale=0.78,>={Stealth[length=2mm]},font=\small]
  % Panel 1: círculo unidad con v1, v2
  \begin{scope}
    \draw[gray!60] (-1.4,0)--(1.4,0) (0,-1.4)--(0,1.4);
    \draw[thick] (0,0) circle (1);
    \draw[->,thick] (0,0)--(45:1) node[above right=-2pt] {$v_1$};
    \draw[->,thick] (0,0)--(-45:1) node[below right=-2pt] {$v_2$};
    \node at (0,-1.9) {$x\in\R^{n}$, $\norm{x}=1$};
  \end{scope}
  \draw[->] (1.8,0)--(3.0,0) node[midway,above] {$V\T$};
  % Panel 2: círculo con ejes canónicos
  \begin{scope}[xshift=4.6cm]
    \draw[gray!60] (-1.4,0)--(1.4,0) (0,-1.4)--(0,1.4);
    \draw[thick] (0,0) circle (1);
    \draw[->,thick] (0,0)--(1,0) node[above right=-2pt] {$e_1$};
    \draw[->,thick] (0,0)--(0,1) node[above right=-2pt] {$e_2$};
    \node at (0,-1.9) {girar};
  \end{scope}
  \draw[->] (6.4,0)--(7.6,0) node[midway,above] {$\Sigma$};
  % Panel 3: elipse alineada con los ejes
  \begin{scope}[xshift=9.6cm]
    \draw[gray!60] (-1.8,0)--(1.8,0) (0,-1.4)--(0,1.4);
    \draw[thick] (0,0) ellipse (1.6 and 0.6);
    \draw[->,thick] (0,0)--(1.6,0) node[above right=-2pt] {$\sigma_1e_1$};
    \draw[->,thick] (0,0)--(0,0.6) node[above left=-2pt] {$\sigma_2e_2$};
    \node at (0,-1.9) {estirar};
  \end{scope}
  \draw[->] (11.6,0)--(12.8,0) node[midway,above] {$U$};
  % Panel 4: elipse girada
  \begin{scope}[xshift=14.6cm]
    \draw[gray!60] (-1.8,0)--(1.8,0) (0,-1.4)--(0,1.4);
    \draw[thick,rotate=30] (0,0) ellipse (1.6 and 0.6);
    \draw[->,thick] (0,0)--(30:1.6) node[above right=-2pt] {$\sigma_1u_1$};
    \draw[->,thick] (0,0)--(120:0.6) node[above left=-2pt] {$\sigma_2u_2$};
    \node at (0,-1.9) {girar: $Ax$};
  \end{scope}
\end{tikzpicture}
\caption{Interpretación geométrica de $A=U\Sigma V\T$ para $m=n=2$. La imagen de la circunferencia unidad es una elipse cuyos semiejes miden $\sigma_1$ y $\sigma_2$ y apuntan en las direcciones $u_1$ y $u_2$.}
\label{fig:svd}
\end{figure}

\begin{observacion}[SVD reducida]
Como las columnas $u_i$ y $v_i$ con $i>r$ se multiplican por ceros, también se cumple $A=U_r\Sigma_rV_r\T$, con $U_r\in\R^{m\times r}$, $\Sigma_r=\diag(\sigma_1,\dots,\sigma_r)$ y $V_r\in\R^{n\times r}$. Es la versión que se almacena en la práctica (\texttt{np.linalg.svd} con \texttt{full\_matrices=False}).
\end{observacion}

\begin{corolario}\label{cor:svd-props}
Sea $A=U\Sigma V\T$ la SVD de $A\in\R^{m\times n}$.
\begin{enumerate}
  \item $\rango(A)$ es el número de valores singulares no nulos.
  \item $\sigma_i(A)=\sqrt{\lambda_i(A\T A)}$, y $AA\T$ tiene los mismos autovalores no nulos que $A\T A$, con autovectores $u_i$.
  \item $\{u_1,\dots,u_r\}$ es una base ortonormal de $\Col(A)$ y $\{v_{r+1},\dots,v_n\}$, una base ortonormal de $\Nuc(A)$.
  \item Si $A$ es simétrica con autovalores $\lambda_i$, sus valores singulares son los $\abs{\lambda_i}$. Si además es semidefinida positiva, la SVD coincide con la descomposición espectral.
\end{enumerate}
\end{corolario}

\begin{proof}
(i) y (iii) se leen en la demostración del \cref{teo:svd}: $Av_i=\sigma_iu_i$ para $i\le r$ y $Av_i=0$ para $i>r$, y $\dim\Nuc(A)=n-r$. (ii) $AA\T=U\Sigma\Sigma\T U\T$ es una descomposición espectral con autovalores $\sigma_i^{2}$; la primera igualdad es la definición de $\sigma_i$. (iv) Si $A=Q\Lambda Q\T$, se toma $V=Q$ y $u_i=\operatorname{signo}(\lambda_i)\,q_i$ (con $u_i=q_i$ si $\lambda_i=0$), de modo que $A=U\abs{\Lambda}V\T$.
\end{proof}

\section{Cálculo a mano}

El procedimiento de la demostración sirve también para calcular la SVD de matrices pequeñas: (1) se forma $A\T A$ o $AA\T$, la que sea de menor orden; (2) se calculan sus autovalores $\sigma_i^{2}$ y autovectores ortonormales; (3) se obtienen los vectores singulares del otro lado con $u_i=Av_i/\sigma_i$ o $v_i=A\T u_i/\sigma_i$.

\begin{ejemplo}\label{ej:svd23}
Sea $A=\begin{pmatrix}3&2&2\\2&3&-2\end{pmatrix}$, la matriz del ejercicio~4.8 de~\cite{mml}.
\begin{enumerate}[label=\arabic*.]
  \item Es más cómodo trabajar con $AA\T=\begin{pmatrix}17&8\\8&17\end{pmatrix}$, de orden~2. Sus autovalores cumplen $(17-\lambda)^{2}=64$, luego $\lambda=25$ y $\lambda=9$, y $\sigma_1=5$, $\sigma_2=3$.
  \item Como en el \cref{ej:S}, los autovectores unitarios son $u_1=(1,1)/\sqrt{2}$ y $u_2=(1,-1)/\sqrt{2}$.
  \item $v_1=A\T u_1/\sigma_1=(5,5,0)/(5\sqrt{2})=(1,1,0)/\sqrt{2}$ y $v_2=A\T u_2/\sigma_2=(1,-1,4)/(3\sqrt{2})$.
\end{enumerate}
Se comprueba que $\norm{v_2}^{2}=(1+1+16)/18=1$ y $v_1\T v_2=(1-1+0)/6=0$. El tercer vector singular derecho, $v_3$, genera $\Nuc(A)$: resolviendo $Av_3=0$ se obtiene $v_3=(2,-2,-1)/3$. La matriz tiene rango~2 y
\[
A=5\,u_1v_1\T+3\,u_2v_2\T.
\]
\end{ejemplo}

\section{Normas matriciales}

\begin{definicion}[Normas de Frobenius y espectral]
Para $A\in\R^{m\times n}$ se definen la \emph{norma de Frobenius} y la \emph{norma espectral}
\[
\norm{A}_F=\Bigl(\sum_{i=1}^{m}\sum_{j=1}^{n}a_{ij}^{2}\Bigr)^{1/2},\qquad
\norm{A}_2=\max_{x\ne 0}\frac{\norm{Ax}}{\norm{x}}.
\]
\end{definicion}

\begin{proposicion}\label{prop:normas-svd}
Para toda $A\in\R^{m\times n}$ con valores singulares $\sigma_1\ge\dots\ge\sigma_r>0$,
\[
\norm{A}_F^{2}=\traza(A\T A)=\sum_{i=1}^{r}\sigma_i^{2},\qquad \norm{A}_2=\sigma_1.
\]
En particular, $\norm{A}_2\le\norm{A}_F\le\sqrt{r}\,\norm{A}_2$.
\end{proposicion}

Las dos igualdades se demuestran en los \cref{eje:4-frobenius,eje:4-espectral}; la segunda es el teorema~4.24 de~\cite{mml}. Ambas normas son invariantes por transformaciones ortogonales: $\norm{Q_1AQ_2}=\norm{A}$ para $Q_1$, $Q_2$ ortogonales, porque los valores singulares no cambian. Las desigualdades finales se siguen de $\sigma_1^{2}\le\sum_i\sigma_i^{2}\le r\sigma_1^{2}$.

\begin{ejemplo}
Para la matriz del \cref{ej:svd23}, $\norm{A}_F^{2}=9+4+4+4+9+4=34=25+9$ y $\norm{A}_2=5$: ningún vector unitario se alarga más de 5 veces, y $v_1$ alcanza ese máximo.
\end{ejemplo}

\begin{observacion}[Conexión con los LLM]
Los valores singulares de una matriz de pesos describen cuánto puede amplificar una capa cada dirección de entrada. Una norma espectral grande indica que pequeñas perturbaciones de la entrada pueden producir grandes cambios en la salida, y el espectro de los pesos se usa para diagnosticar entrenamientos (módulo~11). El \cref{tema:5} usa la SVD para la propiedad más importante en la práctica: truncar la suma $\sum_i\sigma_iu_iv_i\T$ da la mejor aproximación de rango bajo.
\end{observacion}

\begin{lecturas}
\item Strang~\cite{strang-ila}, §7.1 «Singular Values and Singular Vectors».
\item MIT 18.06~\cite{mit1806}, clase~29 «Singular value decomposition».
\item MIT 18.065~\cite{mit18065}, clases~6 «Singular Value Decomposition (SVD)» y 8 «Norms of Vectors and Matrices»; lectura: Strang~\cite{strang-lafd}, §I.8.
\item Deisenroth, Faisal y Ong~\cite{mml}, §4.5 «Singular Value Decomposition» (teorema~4.22) y §4.6 hasta el teorema~4.24.
\item Goodfellow, Bengio y Courville~\cite{gbc}, §2.8 «Singular Value Decomposition».
\end{lecturas}

\section*{Ejercicios del tema 4}
\addcontentsline{toc}{section}{Ejercicios del tema 4}

\begin{ejercicio}{1}{eje:4-vf}
Justifique si son verdaderas o falsas: (a) la SVD solo existe para matrices cuadradas; (b) los valores singulares de una matriz pueden ser negativos; (c) si $A$ es simétrica, sus valores singulares coinciden con sus autovalores; (d) $A$ y $A\T$ tienen los mismos valores singulares.
\end{ejercicio}

\begin{ejercicio}{2}{eje:4-u1}
Sea $A=\begin{pmatrix}1&1\\1&0\end{pmatrix}$. (a) Calcule $A\T A$ y $AA\T$. (b) Calcule sus autovalores y deduzca los valores singulares de $A$. (c) Compruebe que $\sigma_1=(1+\sqrt{5})/2$ y $\sigma_2=(\sqrt{5}-1)/2$, y explique por qué coinciden con los valores absolutos de los autovalores de $A$.
\fuente{adaptado de MIT 18.06SC~\cite{mit1806sc}, sesión «Singular Value Decomposition», problema~29.1.}
\end{ejercicio}

\begin{ejercicio}{2}{eje:4-u5}
Sea $A$ una matriz cuyas columnas $w_1,\dots,w_n$ son ortogonales y no nulas, con $\norm{w_j}=s_j$. Calcule $A\T A$ y deduzca $U$, $\Sigma$ y $V$. Aplíquelo a $A=\begin{pmatrix}2&1\\1&0\\2&-1\end{pmatrix}$.
\fuente{adaptado de MIT 18.06SC~\cite{mit1806sc}, sesión «Singular Value Decomposition», problema~29.2 (caso numérico añadido).}
\end{ejercicio}

\begin{ejercicio}{2}{eje:4-mml}
Calcule la descomposición en valores singulares de $A=\begin{pmatrix}2&2\\-1&1\end{pmatrix}$ y compruebe que $U\Sigma V\T=A$.
\fuente{adaptado de Deisenroth, Faisal y Ong~\cite{mml}, ejercicio~4.9.}
\end{ejercicio}

\begin{ejercicio}{3}{eje:4-frobenius}
Demuestre que $\norm{A}_F^{2}=\traza(A\T A)=\sum_i\sigma_i^{2}$. Indicación: use que $\traza(XY)=\traza(YX)$.
\end{ejercicio}

\begin{ejercicio}{3}{eje:4-autovalores}
Demuestre que, para toda $A\in\R^{m\times n}$, las matrices $A\T A$ y $AA\T$ tienen los mismos autovalores no nulos, con las mismas multiplicidades, sin usar la SVD.
\fuente{adaptado de Deisenroth, Faisal y Ong~\cite{mml}, ejercicio~4.11.}
\end{ejercicio}

\begin{ejercicio}{3}{eje:4-espectral}
Demuestre que $\max_{x\ne 0}\norm{Ax}/\norm{x}=\sigma_1$ y que el máximo se alcanza en $x=v_1$.
\fuente{adaptado de Deisenroth, Faisal y Ong~\cite{mml}, ejercicio~4.12.}
\end{ejercicio}
